\documentclass[11pt,a4paper]{article}
\usepackage[T1]{fontenc}
\usepackage{lmodern}
\usepackage[left=25mm,right=25mm,top=25mm,bottom=25mm,headheight=15pt,headsep=8mm]{geometry}
\usepackage{amsmath,amssymb,amsthm,mathtools}
\usepackage{booktabs,array,microtype,xurl,needspace,etoolbox}
\usepackage[hidelinks]{hyperref}
\usepackage{fancyhdr}
\pagestyle{fancy}
\fancyhf{}
\fancyhead[L]{\small Legendre polynomial irreducibility}
\fancyhead[R]{\small Zikang Deng}
\fancyfoot[C]{\small\thepage}
\renewcommand{\headrulewidth}{0pt}
\setlength{\parindent}{1.5em}
\setlength{\parskip}{0pt plus 0.8pt}
\linespread{1.06}
\setlength{\emergencystretch}{2em}
\allowdisplaybreaks[2]
\numberwithin{equation}{section}
\newtheorem{theorem}{Theorem}[section]
\newtheorem{lemma}[theorem]{Lemma}
\newtheorem{proposition}[theorem]{Proposition}
\theoremstyle{definition}
\newtheorem{definition}[theorem]{Definition}
\pretocmd{\subsection}{\Needspace{5\baselineskip}}{}{}
\pretocmd{\section}{\Needspace{7\baselineskip}}{}{}
\date{}
\hypersetup{pdftitle={A proof of the irreducibility conjecture for Legendre polynomials},pdfauthor={Zikang Deng},pdfsubject={Legendre polynomials and irreducibility}}
\title{\LARGE\bfseries A proof of the irreducibility conjecture for Legendre polynomials}
\author{Zikang Deng\\[4pt]\normalsize Beijing Normal University}
\begin{document}
\maketitle
\begin{abstract}
We prove the irreducibility conjecture for Legendre polynomials proposed by Stieltjes in his 1890 letter to Hermite. Assuming a nontrivial factorization, we take the differences between roots belonging to distinct factors, multiply these differences, and include the leading coefficients to obtain a nonzero integer, the resultant of the two factors. We first bound the exponent with which each odd prime divides this integer, obtaining an upper bound for its odd part. We then use an identity for derivatives at roots, derived from the Legendre differential equation, to prove that the same odd part exceeds another explicit quantity. Comparing the two bounds rules out all factorizations for original degrees $n \ge 100000$. The remaining degrees are covered by Groth's previously established finite-range result.
\end{abstract}

\section{Introduction}\label{sec:intro}

The Legendre polynomials are defined by Rodrigues' formula
\begin{equation}\label{eq:rodrigues}
P_n(x)=\frac{1}{2^n n!}\frac{d^n}{dx^n}(x^2-1)^n,\qquad n\ge0,
\end{equation}
and satisfy \(P_n(1)=1\). Every Legendre polynomial of odd degree has the factor \(x\). Stieltjes' irreducibility conjecture asserts that, after this unavoidable factor has been removed, the Legendre polynomial is irreducible over the rational numbers.

Stieltjes posed this problem in his letter no.~275 to Hermite, dated October~2, 1890 \cite{stieltjes,cullinan}. Wahab \cite{wahab} determined the $2$-adic Newton polygon of the shifted Legendre polynomial and proved irreducibility for several classes of degrees. Cullinan and Hajir \cite{cullinan} further studied its Galois group. In his 2025 doctoral dissertation, Groth \cite{groth} gave an algorithm for testing irreducibility using Newton polygons and other methods, and reported irreducibility for all original degrees \(2\le n\le10^7\); both this range and the convention for even and odd degrees are stated in the abstract of his dissertation. We use this result for the finite range. Our main task is to give a uniform proof for all sufficiently large degrees.

Our main result is as follows.

\begin{theorem}[Stieltjes' irreducibility conjecture]\label{thm:main}
For every integer \(j\ge1\), the polynomials \(P_{2j}(x)\) and \(P_{2j+1}(x)/x\) are irreducible in \(\mathbb Q[x]\). Equivalently, for every integer \(n\ge2\), if \(\varepsilon\in\{0,1\}\) and \(n\equiv\varepsilon\pmod2\), then \(P_n(x)/x^\varepsilon\) is irreducible in \(\mathbb Q[x]\).
\end{theorem}

The proof is by contradiction. Suppose that, after the unavoidable zero root has been removed, the polynomial still factors into two rational factors of smaller degree. We seek a contradiction that can be checked from this assumption. To this end, we normalize each factor to have integer coefficients and consider their resultant. The resultant is a product of differences between roots belonging to different factors, together with the appropriate leading coefficients; it is a nonzero integer.

An integer has a prime factorization. Thus, to obtain an upper bound for the absolute value of the resultant, we may consider the primes one at a time and bound the exponent of each prime in it. This exponent can in turn be computed from the root differences: the higher the power of a prime dividing the difference of two roots, the more closely the roots agree in congruences modulo powers of that prime. For algebraic roots, this statement is made precise by the \(p\)-adic valuation. We shall explain how the derivative of a polynomial at a root can be used to bound the sum of the valuations of these root differences.

On the other hand, all roots of a Legendre polynomial are real. An identity obtained from the differential equation shows that the derivative at each root cannot be too small. Multiplying these derivatives and then dividing out the product of root differences within the same factor gives a lower bound for the resultant. This division requires a determinant estimate, and we shall explain why Chebyshev polynomials are chosen for this purpose.

The exponent of the prime \(2\) can first be computed from the $2$-adic Newton polygon. Thus the two estimates actually concern the same positive integer obtained by removing all factors of \(2\). When the original degree is sufficiently large, the lower bound obtained from the real roots exceeds the upper bound obtained from the prime factorization, making the assumed factorization impossible.

Section~\ref{sec:prelim} introduces the notation and preliminaries. Section~\ref{sec:construction} constructs the resultant and explains why we take its odd part. Section~\ref{sec:upper} first proves local estimates and then sums them to obtain the upper bound. Section~\ref{sec:lower} obtains the lower bound from the real roots. Finally, Section~\ref{sec:conclusion} compares the two bounds directly and proves irreducibility for every \(n\ge100000\); the range \(2\le n<100000\) is covered by Groth's result.

\section{Preliminaries and notation}\label{sec:prelim}

\subsection{Degrees, normalization, and basic notation}

Throughout, fix an integer \(n\ge2\). To treat even and odd degrees together, write
\[
n=2m+\varepsilon,\qquad m=\lfloor n/2\rfloor,\qquad\varepsilon\in\{0,1\}.
\]
Thus no factor is removed when \(\varepsilon=0\), whereas the factor \(x\) is removed when \(\varepsilon=1\). Let \(s(j)=s_2(j)\) denote the number of digits equal to \(1\) in the binary expansion of a nonnegative integer \(j\), with \(s(0)=0\). For example, \(5=(101)_2\), so \(s(5)=2\). Define
\begin{equation}
\kappa_n=2^{n-s(n)},\qquad L_n(x)=\kappa_n P_n(x),\qquad
M_n(x)=\frac{L_n(x)}{x^\varepsilon}.
\label{eq:old1.1}
\end{equation}
Multiplication by a nonzero rational number does not change irreducibility over the rational numbers, so it suffices to prove that \(M_n\) is irreducible. The purpose of this normalization is to obtain a primitive polynomial with integer coefficients and odd leading coefficient. A polynomial is primitive if the greatest common divisor of all its coefficients is \(1\).

For a prime \(p\), let \(v_p\) denote the valuation normalized by \(v_p(p)=1\). For a nonzero integer, \(v_p\) is the exponent of the prime factor \(p\); for a rational number, it is the exponent in the numerator minus the exponent in the denominator. For example, \(v_3(27)=3\) and \(v_3(2/9)=-2\). Set \(v_p(0)=+\infty\). The field \(\mathbb Q_p\) is the field of \(p\)-adic numbers, and \(\overline{\mathbb Q}_p\) is its algebraic closure. In what follows, we extend \(v_p\) to this algebraic closure so that it can be applied to the algebraic roots of polynomials.

We shall also use \(v_p(\alpha-\beta)\) to describe how close two roots are. For example, \(v_3(10-1)=2\), whereas \(v_3(28-1)=3\): the first difference is divisible by \(9\), and the second is divisible by \(27\). The larger the valuation of the difference, the closer the two elements are in the \(p\)-adic sense. The valuation of an algebraic root may be fractional, but it obeys the same rules for addition and multiplication. This notion of closeness differs from distance in the real numbers.

We also explain the meaning of reduction. Place the roots in a finite extension of \(\mathbb Q_p\). The elements of nonnegative valuation form its ring of integers, and the elements of positive valuation form the maximal ideal of this ring. The map from the ring of integers to the field obtained by quotienting out this maximal ideal is called reduction; for ordinary integers, it is reduction modulo \(p\). Two roots of nonnegative valuation have the same reduction if and only if their difference has positive valuation. A root of the reduction means a root of the reduced polynomial in the algebraic closure of this field.

Write \(\operatorname{lc}(f)\) for the leading coefficient of a polynomial \(f\). Write
\[
h_j=\sum_{r=1}^{j}\frac1r\qquad(j\ge1)
\]
for the harmonic sum. All logarithms without a specified base are natural logarithms, and \(\log_p y=\log y/\log p\). The notation used repeatedly below is summarized as follows.

\begin{center}
\begin{tabular}{@{}p{0.17\textwidth}p{0.77\textwidth}@{}}
\toprule
Notation & Meaning\\
\midrule
\(n,\ m,\ \varepsilon\) & Original degree, \(\lfloor n/2\rfloor\), and parity of the degree\\
\(s(j)\) & Binary digit sum of \(j\)\\
\(\kappa_n\) & Power of two \(2^{n-s(n)}\) used to clear denominators\\
\(L_n,\ M_n\) & Integer normalization, and the polynomial with any zero root removed\\
\(v_p,\ h_j\) & Normalized \(p\)-adic valuation and harmonic sum\\
\bottomrule
\end{tabular}
\end{center}

\begin{lemma}\label{lem:primitive}
The polynomials \(L_n,M_n\) are primitive polynomials with integer coefficients and odd leading coefficients.
\end{lemma}

\begin{proof}
The factorial expansion

\begin{equation}
P_n(x)=2^{-n}\sum_{j=0}^{\lfloor n/2\rfloor}
(-1)^j\binom nj\binom{2n-2j}{n}x^{n-2j},
\label{eq:old1.2}
\end{equation}

shows that \(P_n\in\mathbb Z[1/2][x]\). The $2$-adic valuation of the coefficient of the \(j\)-th term is

\[
s(j)+s(n-2j)-n.
\]

Since \(n=2j+(n-2j)\), the digit-sum inequality gives \(s(n)\le s(j)+s(n-2j)\). Thus \(L_n,M_n\in\mathbb Z[x]\), and their leading coefficient

\[
\operatorname{lc}(M_n)=2^{-s(n)}\binom{2n}{n}
\]

is odd. Also, \(M_n(1)=\kappa_n\) is a power of $2$, so \(M_n\) is primitive.
\end{proof}

For example, \(P_5(x)=(63x^5-70x^3+15x)/8\), and hence
\[
M_5(x)=8P_5(x)/x=63x^4-70x^2+15.
\]
In general, \(M_n\) is an even polynomial of degree \(2m\). Therefore there is a unique \(H\in\mathbb Z[t]\) such that
\[
M_n(x)=H(x^2),\qquad\deg H=m.
\]
Here \(t=x^2\) is the squared variable. In the example above, the corresponding polynomial is \(H(t)=63t^2-70t+15\).

\subsection{Real roots and identities for Legendre polynomials}

We record three identities. The first is the Legendre differential equation, which will be used to constrain the derivatives at the roots. The second relates polynomials of consecutive degrees. The third expresses the derivative as a combination of two consecutive polynomials. The last two identities will be used in coefficient calculations after reduction modulo a prime.\begin{equation}
(1-x^2)P_n''-2xP_n'+n(n+1)P_n=0,
\label{eq:old1.3}
\end{equation}

\begin{equation}
(n+1)P_{n+1}=(2n+1)xP_n-nP_{n-1},
\label{eq:old1.4}
\end{equation}

\begin{equation}
(1-x^2)P_n'=n(P_{n-1}-xP_n).
\label{eq:old1.5}
\end{equation}

Integration by parts in Rodrigues' formula shows that \(P_n\) is orthogonal to every polynomial of degree less than \(n\). It follows that all roots of \(P_n\) are simple and lie in \((-1,1)\): otherwise, form a polynomial of degree less than \(n\) from its sign-changing roots in this interval. The product of this polynomial with \(P_n\) has a fixed sign, contradicting orthogonality. The parity identity is \(P_n(-x)=(-1)^nP_n(x)\), so, apart from the zero root when the degree is odd, the roots occur in pairs \(\pm x\).

We shall also use

\begin{equation}
|P_j(x)|\le1\quad(-1\le x\le1).
\label{eq:old1.6}
\end{equation}

For example, the Laplace representation obtained from the generating function,

\[
P_j(x)=\frac1\pi\int_0^\pi
\bigl(x+i\sqrt{1-x^2}\cos\theta\bigr)^j\,d\theta
\]

immediately proves \eqref{eq:old1.6}, because the complex number in parentheses has modulus at most $1$. These formulas may also be found in \cite{dlmf}.

In particular, all roots of \(H\) are simple and lie in \((0,1)\). Their location will be used for the lower bound for the resultant in Section~\ref{sec:lower}.

\subsection{Newton polygons and the binary structure}

To determine which roots can belong to a rational factor, we first factor the polynomial over \(\mathbb Q_p\). The Newton polygon is drawn using the \(p\)-adic valuations of the coefficients, and each of its edges gives a group of roots with the same valuation. We mainly use edges whose vertical change is one: the following lemma shows that such a group of roots forms an irreducible factor over \(\mathbb Q_p\). Consequently, any rational factor must contain either this entire group of roots or none of them.

\begin{definition}\label{def:np}
Let \(f(X)=\sum_{j=0}^d a_jX^j\in\mathbb Q_p[X]\), with \(a_d\ne0\). The \(p\)-adic Newton polygon of \(f\) is the lower convex hull of the points \((j,v_p(a_j))\) for which \(a_j\ne0\). The horizontal length of an edge is called its width.
\end{definition}

\begin{lemma}\label{lem:np}
An edge of the Newton polygon with slope \(-\lambda\) and width \(q\) corresponds to exactly \(q\) roots of valuation \(\lambda\), counted with multiplicity. If the vertical change along the edge is \(1\) or \(-1\), these roots form an irreducible factor over \(\mathbb Q_p\).
\end{lemma}

\begin{proof}
The first assertion is the root-valuation theorem for Newton polygons. For the second assertion, the roots corresponding to the edge have valuation \(\lambda=\pm1/q\). The conjugate roots of an irreducible factor over \(\mathbb Q_p\) have the same valuation. If its degree is \(d\), then, after making it monic, the valuation of its constant term is \(d\lambda\in\mathbb Z\), so \(q\mid d\). Since there are only \(q\) roots of this valuation in total, the edge can correspond to only one irreducible factor.
\end{proof}

For Legendre polynomials, Wahab's theorem relates these edges directly to the binary expansion of the degree \cite{wahab,cullinan}. For use below, we give a direct proof.

\begin{theorem}[Wahab]\label{thm:wahab}
Let

\[
Q_n(X)=P_n(2X+1)
=\sum_{j=0}^n c_jX^j,\qquad
c_j=\frac{(n+j)!}{(n-j)!(j!)^2}.
\]

Write \(n\) as a sum of distinct powers of two,

\[
n=q_1+\cdots+q_s,\qquad q_1>\cdots>q_s.
\]

Then the successive edges of the $2$-adic Newton polygon of \(Q_n\) have widths \(q_i\) and height $1$, and hence slopes \(1/q_i\).
\end{theorem}

\begin{proof}
Legendre's formula for the valuation of a factorial gives

\begin{equation}
v_2(c_j)=s(n-j)+2s(j)-s(n+j)\ge s(j),
\label{eq:old2.1}
\end{equation}

where the inequality follows from \(n+j=(n-j)+2j\). Put \(K_t=q_1+\cdots+q_t\) and \(K_0=0\). Since the sets of nonzero positions in the binary expansions of \(2K_t\) and \(n-K_t\) are disjoint,

\[
s(n+K_t)=t+s(n-K_t),\qquad v_2(c_{K_t})=t.
\]

When \(K_t<j<K_{t+1}\), write \(j=K_t+u\), where \(0<u<q_{t+1}\). Then \(s(j)=t+s(u)\ge t+1\). Thus these coefficient points lie strictly above the segment joining \((K_t,t)\) and \((K_{t+1},t+1)\), proving the asserted description of the polygon.

The root-valuation theorem for Newton polygons shows that each \(q_i\) corresponds to exactly \(q_i\) roots satisfying \(v_2(X)=-1/q_i\), and that these roots form an irreducible factor over \(\mathbb Q_2\). The last assertion can also be seen directly from the denominator: the roots of any irreducible factor corresponding to this slope have the same valuation, and the valuation of its constant term forces its degree to be divisible by \(q_i\); the total width at this slope is only \(q_i\).
\end{proof}

\subsection{An auxiliary identity connecting the two estimates}

The reason for studying derivatives at roots can be seen directly from a factorization. Suppose that a polynomial has distinct roots \(\alpha_1,\ldots,\alpha_d\) and is written as
\[
f(x)=c\prod_{j=1}^d(x-\alpha_j).
\]
Differentiate and substitute \(x=\alpha_i\). Every term vanishes except the term in which the factor \(x-\alpha_i\) was differentiated. Thus
\[
f'(\alpha_i)=c\prod_{j\ne i}(\alpha_i-\alpha_j).
\]
This shows that an estimate for the derivative gives an estimate for a whole product of root differences. For Legendre polynomials, the required identity is obtained from the following explicit auxiliary polynomial:
\begin{equation}
W_{n-1}(x)=\sum_{j=1}^n\frac1jP_{j-1}(x)P_{n-j}(x).
\label{eq:old3.1}
\end{equation}
The subscript \(n-1\) gives its degree. This definition has two properties that make estimates possible. First, for an odd prime \(p\), the denominators of the coefficients of \(P_j\) contain no factor of \(p\), so the factors \(1/j\) in the sum are the only possible source of denominators divisible by \(p\). Second, on the real interval \([-1,1]\), the absolute value of each \(P_j\) is at most \(1\), so the absolute value of the whole sum is at most \(1+1/2+\cdots+1/n\). The following identity transfers these two properties to the derivatives at the roots.

\begin{lemma}[Wronskian identity]\label{lem:wronskian}
For every integer \(n\ge1\), we have

\begin{equation}
(1-x^2)(P_n'W_{n-1}-P_nW_{n-1}')=1-P_n^2,
\label{eq:old3.2}
\end{equation}

\begin{equation}
|W_{n-1}(x)|\le h_n\quad(-1\le x\le1).
\label{eq:old3.3}
\end{equation}
\end{lemma}

\begin{proof}
Equation~\eqref{eq:old3.3} follows directly from \eqref{eq:old1.6}. We now prove \eqref{eq:old3.2}.

Let \(\mathcal F(x,T)=\sum_{n\ge0}P_n(x)T^n=(1-2xT+T^2)^{-1/2}\). By \eqref{eq:old3.1},

\[
\sum_{n\ge1}W_{n-1}(x)T^n
=\mathcal F(x,T)\int_0^T\mathcal F(x,u)\,du.
\]

Denote the right-hand side by \(G(x,T)\). It satisfies \((1-2xT+T^2)\partial_T G-(x-T)G=1\). Comparing coefficients gives \(W_{-1}=0,W_0=1\), and shows that \(W\) satisfies the same recurrence as \eqref{eq:old1.4}.

For \(z\notin[-1,1]\), define

\begin{equation}
\mathsf Q_n(z)=\frac12\int_{-1}^1\frac{P_n(t)}{z-t}\,dt.
\label{eq:old3.4}
\end{equation}

Orthogonality and \eqref{eq:old1.4} give the same recurrence and the initial values, and hence

\begin{equation}
\mathsf Q_n(z)=\frac12P_n(z)\log\frac{z+1}{z-1}-W_{n-1}(z).
\label{eq:old3.5}
\end{equation}

The kernel \(1/(z-t)\) satisfies \(\mathcal L_zK=\mathcal L_tK\), where \(\mathcal L=(1-x^2)\partial_x^2-2x\partial_x\). Integrating by parts in \eqref{eq:old3.4}, the boundary terms vanish because \(1-t^2=0\) at the endpoints, so \(\mathsf Q_n\) also satisfies the Legendre equation. Therefore

\[
(1-z^2)(P_n\mathsf Q_n'-P_n'\mathsf Q_n)
\]

is constant. Substituting \eqref{eq:old3.5}, this expression becomes

\[
P_n^2+(1-z^2)(P_n'W_{n-1}-P_nW_{n-1}'),
\]

which is a polynomial taking the value $1$ at \(z=1\). Thus the constant is $1$, proving \eqref{eq:old3.2}.
\end{proof}

In particular, when \(P_n(\alpha)=0\), the identity becomes
\begin{equation}\label{eq:root_identity}
(1-\alpha^2)P_n'(\alpha)W_{n-1}(\alpha)=1.
\end{equation}
Over the real numbers, the right-hand side is \(1\), while \(|W_{n-1}(\alpha)|\) has an upper bound, so \(|P_n'(\alpha)|\) cannot be arbitrarily small. At an odd prime, taking valuations on both sides shows that the valuations of the three factors sum to \(0\). Thus lower bounds for the valuations of the other two factors give an upper bound for the valuation of the derivative. Sections~\ref{sec:upper} and~\ref{sec:lower} carry out these two estimates, respectively.

\section{Construction of the quantity to be compared}\label{sec:construction}

From this section onward, suppose that \(M_n\) has a nontrivial factorization over \(\mathbb Q\). Our aim is to construct a nonzero integer to which we can apply both information about the real roots and information about divisibility by primes. The construction has three steps: use the binary structure to determine the form of the factors, take their resultant, and then remove the power of \(2\) already determined by the binary structure.

\subsection{From a rational factorization to the squared variable}

\begin{lemma}\label{lem:even}
Every factor of \(M_n\) over \(\mathbb Q\) is an even polynomial.
\end{lemma}

\begin{proof}
Return to the original variable \(x=2X+1\). The roots associated with \(q_i\) satisfy

\begin{equation}
v_2(1-x)=1-1/q_i.
\label{eq:old2.2}
\end{equation}

If \(q_i\ge2\), the right-hand side is less than \(1\), and hence

\[
v_2(1+x)=v_2\bigl(2-(1-x)\bigr)=1-1/q_i.
\]

Thus the map \(x\mapsto-x\) preserves the set of roots with this valuation. For each \(q_i\), these roots form a single irreducible factor over \(\mathbb Q_2\), so its root set is stable under negation. A rational factor must be a product of some of these irreducible factors over \(\mathbb Q_2\). Consequently, after the possible zero root is removed, every rational factor is an even polynomial. If \(q_i=1\) occurs, its unique root is the zero root of the polynomial of odd degree, which has already been removed in \(M_n\).
\end{proof}

By Lemma~\ref{lem:even}, if we assume below that \(M_n\) is reducible, Gauss's lemma allows us to write

\begin{equation}
M_n(x)=F(x)G(x)=A(x^2)B(x^2),
\label{eq:old2.3}
\end{equation}

where \(F,G,A,B\) are primitive polynomials with integer coefficients and positive leading coefficients. Write

\[
H(t)=A(t)B(t),\qquad M_n(x)=H(x^2).
\]

Also put

\[
\deg A=k,\quad\deg B=e=m-k,\quad1\le k\le e.
\]

The powers of \(2\) occurring in the binary expansion of \(m\) are partitioned into disjoint sets \(S,T\) such that

\begin{equation}
k=\sum_{u\in S}u,\qquad e=\sum_{v\in T}v.
\label{eq:old2.4}
\end{equation}

Thus \(s(m)=s(k)+s(e)\), \(|S|=s(k)\), and \(|T|=s(e)\). All roots of \(A,B\) lie in \((0,1)\).

Thus \(A,B\) are the two nonconstant factors in the assumed factorization in the variable \(t\), and \(k,e\) are their degrees. Interchanging the factors ensures that \(k\le e\). Put \(a=\operatorname{lc}(A)\) and \(b=\operatorname{lc}(B)\); both are positive odd integers. We next determine the endpoint values. These will control the contribution from the leading coefficients in Section~\ref{sec:lower}.

\begin{lemma}\label{lem:endpoint}
With the notation above, \(A(1)=2^{2k-s(k)}\) and \(B(1)=2^{2e-s(e)}\). If \(a=\operatorname{lc}(A)\), then \(a>A(1)\); in particular, \(a\ge3\) when \(k=1\).
\end{lemma}

\begin{proof}
Since \(F(1)G(1)=\kappa_n\) and both factors take positive integer values at \(1\), these values are powers of \(2\). The leading coefficient of \(F\) is odd. Summing \eqref{eq:old2.2} over its roots gives

\begin{equation}
F(1)=A(1)=2^{2k-s(k)},\qquad
G(1)=B(1)=2^{2e-s(e)}.
\label{eq:old2.5}
\end{equation}

If \(a=\operatorname{lc}(A)\), then

\begin{equation}
A(1)=a\prod_{A(t_i)=0}(1-t_i)<a.
\label{eq:old2.6}
\end{equation}

In particular, \(a\) is a positive odd integer. When \(k=1\), we have \(A(1)=2\), so \(a\ge3\).
\end{proof}

\subsection{Why we use the resultant}

Let \(t_1,\ldots,t_k\) be the roots of \(A\), and let \(u_1,\ldots,u_e\) be the roots of \(B\). An individual root difference need not be rational. However, multiplying all root differences between the two factors, together with the appropriate powers of the leading coefficients, gives the resultant
\begin{equation}\label{eq:resultant_product}
R=\operatorname{Res}_t(A,B)
  =a^e b^k\prod_{i=1}^k\prod_{j=1}^e(t_i-u_j).
\end{equation}
The resultant is also the determinant of the Sylvester matrix, whose entries are coefficients of \(A,B\). Therefore \(R\in\mathbb Z\). Since \(H=AB\) has no repeated roots, the two root sets are disjoint, and hence \(R\ne0\).

This is why we choose the resultant: it retains the distances between the two sets of real roots while also being an integer with a prime factorization. Thus
\[
\log|R|=\sum_p v_p(R)\log p.
\]
The left-hand side can be estimated using the real roots, whereas the right-hand side can be estimated one prime at a time.

\subsection{Removing the known power of two}

The prime \(2\) has a special role for Legendre polynomials: it controls the normalization, the endpoint values, and the valuations of the roots. Computing its contribution to \(R\) exactly at the outset allows us to concentrate the subsequent comparison on the odd primes.

\begin{lemma}\label{lem:two}
For the factorization above, the resultant \(R\) satisfies \eqref{eq:old2.7}--\eqref{eq:old2.9}.
\end{lemma}

\begin{proof}
For distinct powers \(u,v\) in the binary expansion, the corresponding squared variables \(t=x^2\) satisfy

\[
v_2(1-t)=2-1/u,
\qquad
v_2(t-t')=2-1/\min(u,v).
\]

Since both leading coefficients are odd, we obtain

\begin{equation}
v_2(R)=\sum_{u\in S,v\in T}\bigl(2uv-\max(u,v)\bigr)
=2ke-\sum_{u\in S,v\in T}\max(u,v).
\label{eq:old2.7}
\end{equation}

In particular, since \(\max(u,v)\ge v\),

\begin{equation}
v_2(R)\le(2k-s(k))e.
\label{eq:old2.8}
\end{equation}

Taking the product over the roots directly also gives

\begin{equation}
\operatorname{Res}_x(A(x^2),B(x^2))=R^2.
\label{eq:old2.9}
\end{equation}
\end{proof}

Define the odd part of the resultant by
\begin{equation}\label{eq:odd_part}
R_{\mathrm{odd}}=\frac{|R|}{2^{v_2(R)}}.
\end{equation}
This is a positive odd integer, and
\begin{equation}\label{eq:odd_log_sum}
\log R_{\mathrm{odd}}=\sum_{p\text{ an odd prime}}v_p(R)\log p.
\end{equation}
From now on, all upper and lower bounds concern this same number. Equation~\eqref{eq:old2.9} also shows that, after estimating the valuation of the resultant of the two factors in the original variable \(x\), we must divide by \(2\) to obtain \(v_p(R)\) in the squared variable.

\section{An upper bound for the odd part of the resultant}\label{sec:upper}

In this section we seek an upper bound for the positive integer \(R_{\mathrm{odd}}\). Equation~\eqref{eq:odd_log_sum} expresses its logarithm as a sum of terms \(v_p(R)\log p\). We therefore first fix an odd prime \(p\), prove that \(v_p(R)\) cannot be too large, and then sum over all odd primes.

Why do root differences determine this exponent? Taking valuations in \eqref{eq:resultant_product} gives
\[
v_p(R)=e\,v_p(a)+k\,v_p(b)
       +\sum_{i=1}^{k}\sum_{j=1}^{e}v_p(t_i-u_j).
\]
Here \(a,b\) are the leading coefficients of the two factors, and \(t_i,u_j\) are their respective roots. If neither leading coefficient is divisible by \(p\), the first two terms vanish, and the exponent is exactly the sum of these root-difference valuations. In general, the first two terms must also be included; the proof below will account for them. The total upper bound that we shall prove is the following.

\begin{proposition}\label{prop:upper}
Under the assumed factorization of Section~\ref{sec:construction}, we have
\begin{equation}
\log R_{\mathrm{odd}}
<k\sqrt{2n}\log(2n)+2m\log2\,h_{2k-1}.
\label{eq:old7.3}
\end{equation}
\end{proposition}

To obtain this upper bound, we first establish two facts. First, for every odd prime, we have \(v_p(R)<k\log_p(2n)\); this follows from the identity \eqref{eq:root_identity} for the derivative at a root. Second, when \(p^2>2n\), we also have \(v_p(R)\le\lfloor m/p\rfloor\). This requires a closer examination of which roots have the same reduction and whether such roots can belong to different factors.

The first bound limits the exponent of each prime, but we also need the second bound to constrain the large primes. For example, if \(p>m\) and the condition for the second bound holds, then \(\lfloor m/p\rfloor=0\), so this prime does not divide \(R\) at all. Combining the two estimates will give the total upper bound above.

\subsection{A uniform estimate at all odd primes}

Fix an odd prime \(p\). We divide the proof into two cases according to the valuation of a root \(\alpha\).

If \(v_p(\alpha)\ge0\), substituting \(\alpha\) into \(P_j\) introduces no additional powers of \(p\) in the denominator. We can therefore estimate \(W_{n-1}(\alpha)\) directly from the definition of \(W_{n-1}\), and then estimate the derivative using \eqref{eq:root_identity}.

If \(v_p(\alpha)<0\), direct substitution makes higher powers introduce larger denominators, so the first method does not apply directly. Instead, we use \(\beta=1/\alpha\), so that \(v_p(\beta)>0\), and rewrite the polynomials accordingly as polynomials in the reciprocal variable. After obtaining upper bounds for the derivative valuations in the two cases, we explain how they combine to give an upper bound for \(v_p(R)\).

\begin{lemma}\label{prop:uniform}
For every factorization \eqref{eq:old2.3} and every odd prime \(p\),

\begin{equation}
v_p(R)<k\log_p(2n).
\label{eq:old4.1}
\end{equation}
\end{lemma}

\begin{proof}
\textbf{(i) Roots satisfying \(v_p(\alpha)\ge0\).}

Suppose that \(P_n(\alpha)=0\) and \(v_p(\alpha)\ge0\). The coefficients of \(P_n\) are integral at the odd prime \(p\), whereas \(P_n(\pm1)=(\pm1)^n\) has valuation zero. Thus \(\alpha\) cannot reduce to \(1\) or \(-1\): otherwise substitution of \(\alpha\) would also give a value with nonzero reduction, contrary to \(P_n(\alpha)=0\). Hence \(v_p(1-\alpha^2)=0\), so \(1-\alpha^2\) is a unit.

Since \(P_j\in\mathbb Z_p[x]\) and \(v_p(\alpha)\ge0\), each \(P_j(\alpha)\) has nonnegative valuation. In the \(j\)-th summand of \(W_{n-1}(\alpha)\), only the denominator \(j\) can decrease the valuation, and \(v_p(j)\le\lfloor\log_p n\rfloor\). The valuation of a sum is at least the smallest valuation of its summands, so \eqref{eq:old3.1} gives

\[
v_p(W_{n-1}(\alpha))\ge-\lfloor\log_p n\rfloor.
\]

Substituting the root into \eqref{eq:old3.2}, we obtain

\begin{equation}
v_p(P_n'(\alpha))
=-v_p(W_{n-1}(\alpha))
\le\lfloor\log_p n\rfloor<\log_p(2n).
\label{eq:old4.2}
\end{equation}

\textbf{(ii) Roots satisfying \(v_p(\alpha)<0\): passing to reciprocals.}

If \(v_p(\alpha)<0\), put \(\beta=1/\alpha\), \(\rho=v_p(\beta)>0\), and define

\[
\widehat P_n(Y)=Y^nP_n(1/Y),\qquad
\widehat W_n(Y)=Y^{n-1}W_{n-1}(1/Y).
\]

The factors \(Y^n\) and \(Y^{n-1}\) clear the negative powers introduced by substituting \(x=1/Y\), so both expressions with hats are still polynomials. The number \(\alpha\) is a root of \(P_n\) if and only if \(\beta\) is a root of \(\widehat P_n\). Rewriting \eqref{eq:old3.2} in terms of \(Y\) and then substituting this root gives

\begin{equation}
(1-\beta^2)\widehat P_n'(\beta)\widehat W_n(\beta)
=\beta^{2n-1}.
\label{eq:old4.3}
\end{equation}

The right-hand side is now \(\beta^{2n-1}\), whose valuation is already large. To obtain an upper bound for the derivative valuation from this identity, we must prove that the other factor, \(\widehat W_n(\beta)\), also has sufficiently large valuation. The integral representation of the function of the second kind provides precisely this information, because orthogonality makes many of the initial terms of its expansion vanish. More explicitly, expand the kernel as
\[
\frac1{1/Y-t}=\frac{Y}{1-tY}=\sum_{r\ge0}t^rY^{r+1}.
\]
For \(r<n\), orthogonality gives \(\int_{-1}^1 t^rP_n(t)\,dt=0\). When \(r-n\) is odd, the integrand is odd, so the integral also vanishes. Thus the first possibly nonzero term starts at \(Y^{n+1}\), and subsequent terms occur at every other degree. This gives

\begin{equation}
\mathsf Q_n(1/Y)=\sum_{j\ge0}c_jY^{n+2j+1},\qquad
c_j=\frac12\int_{-1}^1t^{n+2j}P_n(t)\,dt.
\label{eq:old4.4}
\end{equation}

These are ordinary real integrals, and the resulting \(c_j\) are rational numbers. Since \(P_n\in\mathbb Z_p[x]\), termwise integration introduces only integer denominators not exceeding \(2n+2j+1\). Hence

\begin{equation}
v_p(c_j)\ge-\lfloor\log_p(2n+2j+1)\rfloor.
\label{eq:old4.5}
\end{equation}

If \(v_p(Y)>0\), the valuation of \(Y^{n+2j+1}\) grows linearly with \(j\), whereas the possible loss in the coefficient valuation grows only logarithmically. Thus the series \eqref{eq:old4.4} converges. Equation~\eqref{eq:old3.5} may first be interpreted as a formal series identity with rational coefficients, in which
\[
\log\frac{1+Y}{1-Y}=2\left(Y+\frac{Y^3}{3}+\cdots\right)
\]
also converges when \(v_p(Y)>0\). Upon substituting \(Y=\beta\), the term containing \(P_n(1/\beta)=P_n(\alpha)=0\) vanishes. Multiplying by \(\beta^{n-1}\) then gives

\[
\widehat W_n(\beta)
=-\sum_{j\ge0}c_j\beta^{2n+2j}.
\]

The valuations of the terms in this series imply
\[
v_p(\widehat W_n(\beta))
\ge\inf_{j\ge0}\bigl\{(2n+2j)\rho-\log_p(2n+2j+1)\bigr\}.
\]
Since \(\rho>0\), the valuation of \(1-\beta^2\) is zero. Taking valuations in \eqref{eq:old4.3} yields
\[
v_p(\widehat P_n'(\beta))
=(2n-1)\rho-v_p(\widehat W_n(\beta)).
\]
Substituting the preceding lower bound gives

\begin{equation}
v_p(\widehat P_n'(\beta))
\le\sup_{j\ge0}\left\{\log_p(2n+2j+1)-(2j+1)\rho\right\}.
\label{eq:old4.6}
\end{equation}

We still need a uniform lower bound for the positive number \(\rho=v_p(\beta)\). The vertices of the Newton polygon of \(\widehat P_n\) have integer coordinates. The absolute vertical change along a nonhorizontal edge is at least \(1\), whereas its horizontal length is at most the degree of the polynomial, which is at most \(n\). The root-valuation theorem therefore gives \(\rho\ge1/n\). Substituting this lower bound and using \(\log(1+u)\le u\), we find that each term in \eqref{eq:old4.6} is at most

\[
\log_p(2n)-\frac{2j+1}{n}
\left(1-\frac1{2\log p}\right).
\]

Since \(p\ge3\), the expression in parentheses is strictly positive. We therefore obtain a positive gap independent of the index:

\begin{equation}
v_p(\widehat P_n'(\beta))
\le\log_p(2n)-\frac1n\left(1-\frac1{2\log p}\right)
<\log_p(2n).
\label{eq:old4.7}
\end{equation}

Thus the same strict upper bound also holds for nonintegral roots arising when the leading coefficient is not a unit.

\textbf{(iii) From derivative estimates to a resultant estimate.}

The first two parts estimate the derivative at a root in the original variable and in the reciprocal variable, respectively. We now express them by a single sum. For distinct roots \(\alpha,\gamma\), define
\begin{equation}
\delta_p(\alpha,\gamma)
=v_p(\alpha-\gamma)-\min(v_p(\alpha),0)-\min(v_p(\gamma),0)\ge0.
\label{eq:old4.8}
\end{equation}
This definition can be understood directly through three cases:
\[
\delta_p(\alpha,\gamma)=
\begin{cases}
v_p(\alpha-\gamma),&v_p(\alpha),v_p(\gamma)\ge0,\\
v_p(\alpha^{-1}-\gamma^{-1}),&v_p(\alpha),v_p(\gamma)<0,\\
0,&\begin{array}{l}\text{one valuation is nonnegative}\\\text{and the other is negative}.\end{array}
\end{cases}
\]
In the first case, we retain the original root difference. In the second, we use the difference of the reciprocals, since
\(\alpha^{-1}-\gamma^{-1}=(\gamma-\alpha)/(\alpha\gamma)\).
The third case uses the basic property of valuations that, when two numbers have different valuations, the valuation of their difference equals the smaller one. Thus all three cases give nonnegative numbers. When one root is zero, we still use the first or third case, with the convention \(v_p(0)=+\infty\).

We next explain the relation between this quantity and the derivative. For each root \(\gamma\), choose the homogeneous linear factor
\[
\ell_\gamma(X,Z)=
\begin{cases}
X-\gamma Z,&v_p(\gamma)\ge0,\\
\gamma^{-1}X-Z,&v_p(\gamma)<0.
\end{cases}
\]
In either expression, all coefficients have nonnegative valuation, and at least one coefficient has valuation zero. Homogenizing \(P_n\) means rewriting it as \(Z^nP_n(X/Z)\). This homogeneous polynomial is the product of the linear factors above times a scalar. Since \(P_n(1)=1\), the polynomial \(P_n\) is primitive at \(p\), meaning that at least one coefficient has valuation zero. The scalar is therefore also a unit. Here we use the multiplicativity of the Gauss norm: if the maximum \(p\)-adic absolute value of the coefficients of each of two polynomials is \(1\), then the maximum absolute value of the coefficients of their product is also \(1\), where \(|z|_p=p^{-v_p(z)}\).

If \(v_p(\alpha)\ge0\), differentiate in the coordinates \((X,Z)=(x,1)\). The factor belonging to \(\alpha\) is \(x-\alpha\). After differentiation and evaluation at \(x=\alpha\), what remains is the product of the values of all the other factors at this root. If \(v_p(\alpha)<0\), differentiate in the coordinates \((X,Z)=(1,Y)\). The factor belonging to \(\alpha\) is \(\alpha^{-1}-Y\), whose derivative is the unit \(-1\), and again all the other factors remain. Consequently,
\[
\sum_{\gamma\ne\alpha}\delta_p(\alpha,\gamma)=
\begin{cases}
v_p(P_n'(\alpha)),&v_p(\alpha)\ge0,\\
v_p(\widehat P_n'(\alpha^{-1})),&v_p(\alpha)<0.
\end{cases}
\]
The two quantities on the right are exactly the derivative valuations estimated in the first two parts.

Now consider the resultant. Since \(F\) and \(G\) are primitive integer polynomials, the same Gauss-norm argument shows that the scalar multiplying the normalized linear factors of each is a unit. The resultant of two linear factors is the two-by-two determinant of their coefficients. Substituting each of the two types of linear factor above shows that its valuation is exactly \(\delta_p(\alpha,\gamma)\). Since resultants are multiplicative with respect to products of polynomials, we obtain
\[
v_p\operatorname{Res}_x(F,G)
=\sum_{\substack{F(\alpha)=0\\G(\gamma)=0}}
\delta_p(\alpha,\gamma).
\]
For a fixed root \(\alpha\) of \(F\), this sum includes only the roots belonging to \(G\). The sum given by the derivative, on the other hand, runs over all the other roots of \(P_n\). Since every summand is nonnegative, the partial sum in the resultant cannot exceed the full sum given by the derivative. Removing the possible zero root of \(P_n\) also removes only a nonnegative term.

Finally, \(F=A(x^2)\) has \(2k\) roots. Applying \eqref{eq:old4.2} or \eqref{eq:old4.7} at each root and summing gives
\[
v_p\operatorname{Res}_x(F,G)<2k\log_p(2n).
\]
Since \(\operatorname{Res}_x(F,G)=R^2\), the left-hand side equals \(2v_p(R)\). Dividing by \(2\) proves \eqref{eq:old4.1}.
\end{proof}

\subsection{A stronger estimate at large primes}

The uniform estimate bounds the exponent of an individual prime, but a stronger bound is needed for large primes. We shall show that roots can contribute to \(v_p(R)\) only when they coincide after reduction. Moreover, when \(p^2>2n\), the group of roots reducing to any one nonzero multiple root contributes at most \(1\) to the valuation of the resultant in the original variable.

Indeed, the difference of two roots with distinct reductions is a unit, and hence makes no contribution to the valuation of the resultant. It therefore remains to consider, one group at a time, roots with the same reduction and to determine whether the roots in such a group can belong to two different rational factors.

\begin{lemma}\label{prop:large}
If an odd prime \(p\) satisfies \(p^2>2n\), then

\begin{equation}
v_p(R)\le\left\lfloor\frac{n}{2p}\right\rfloor
=\left\lfloor\frac mp\right\rfloor.
\label{eq:old5.1}
\end{equation}
\end{lemma}

\begin{proof}
\textbf{(i) Factorization modulo a prime.}

Write \(n=\ell p+r\), where \(0\le r<p\). Then \(\ell<p/2\). The classical congruences are \cite{cullinan,noe}

\begin{equation}
\overline P_{\ell p+r}=\overline P_{\ell}^{\,p}\overline P_r,
\qquad \overline P_r=\overline P_{p-1-r}.
\label{eq:old5.2}
\end{equation}

Both identities can also be proved directly. In characteristic \(p\), the generating function \(\mathcal F=(1-2xT+T^2)^{-1/2}\) satisfies

\[
\mathcal F(x,T)=(1-2xT+T^2)^{(p-1)/2}\mathcal F(x,T)^p
=(1-2xT+T^2)^{(p-1)/2}\mathcal F(x^p,T^p).
\]

The first factor has degree \(p-1\) in \(T\), and its coefficients in these degrees are precisely \(\overline P_r(x)\). Comparing the coefficients of \(T^{\ell p+r}\) gives the first identity. Since \(1-2xT+T^2\) is palindromic in \(T\), its \((p-1)/2\)-th power has equal coefficients in degrees \(r\) and \(p-1-r\). This gives the second identity.

Thus the actual degree of \(\overline P_r\) is \(u=\min(r,p-1-r)\le(p-1)/2\). Both \(\overline P_{\ell}\) and \(\overline P_u\) have only simple roots. The endpoints \(\pm1\) are not roots. If any other point had multiplicity \(2\le d<p\), substituting the lowest-order term into \eqref{eq:old1.3} would force \(d(d-1)=0\pmod p\), a contradiction.

Group together the roots that reduce to the same root of the reduction. By \eqref{eq:old5.2}, every nonzero multiple root of the reduction comes from \(\overline P_{\ell}\). Since both \(\overline P_{\ell}\) and \(\overline P_r\) have only simple roots, its multiplicity in \(\overline P_n\) has only two possibilities: it is \(p\) if the point is not a root of \(\overline P_r\), and it is \(p+1\) if the point is also a root of \(\overline P_r\). When \(\ell=0,1\), there are no nonzero multiple roots of the reduction.

If the roots in such a group form an irreducible polynomial, the rational factors \(F,G\) cannot each contain only part of the group: the entire group must belong to the same factor. In that case, the group produces no root differences between \(F\) and \(G\). It can contribute to the valuation of the resultant only if it can split into two factors.

\textbf{(ii) Zero and infinity.}

Since \(p^2>2n\), only the first term in the factorial valuation formula is needed. Equation \eqref{eq:old1.2} gives

\begin{equation}
v_p(\operatorname{lc}M_n)=\left\lfloor\frac{2r}{p}\right\rfloor\in\{0,1\},
\qquad v_p(M_n(0))=\ell\bmod2\in\{0,1\}.
\label{eq:old5.3}
\end{equation}

For even \(n\), the second valuation is
\(\lfloor n/p\rfloor-2\lfloor n/(2p)\rfloor\); for odd \(n\), it is

\[
\left\lfloor\frac{n+1}{p}\right\rfloor
-\left\lfloor\frac{n-1}{2p}\right\rfloor
-\left\lfloor\frac{n+1}{2p}\right\rfloor.
\]

Substituting \(n=\ell p+r\) reduces both expressions to \(\ell\bmod2\), including the boundary cases \(r=0,p-1\).

Since \(M_n=FG\), we have
\[
M_n(0)=F(0)G(0),\qquad
\operatorname{lc}(M_n)=\operatorname{lc}(F)\operatorname{lc}(G).
\]
The valuations of all factors on the right are nonnegative integers, whereas \eqref{eq:old5.3} shows that the valuations of both products are at most \(1\). Thus the two constant terms cannot both be divisible by \(p\), and neither can the two leading coefficients.

The reduction of a polynomial has a root at \(0\) if and only if its constant term is divisible by \(p\). In the reciprocal variable, the original leading coefficient becomes the constant term. Thus the reduction has a root at infinity if and only if the leading coefficient is divisible by \(p\). It follows that the reductions of \(F,G\) cannot share a root at either of these two points. There is consequently no positive valuation of a root difference to count at either point.

\textbf{(iii) An identity at a nonzero root of the reduction.}

Fix a nonzero root of \(\overline P_{\ell}\). Choose a finite unramified extension \(K/\mathbb Q_p\) over whose residue field \(\overline P_{\ell}\) splits. Since the chosen root of the reduction is simple, Hensel's lemma lifts it to a root \(\alpha\) of \(P_{\ell}\). We shall expand \(P_n\) about this \(\alpha\). It is an exact root of \(P_{\ell}\), but need not be a root of \(P_n\). Our next task is precisely to determine how many times \(p\) divides \(P_n(\alpha)\) and \(P_n'(\alpha)\). Put

\[
u_j=P_{\ell p+j}(\alpha)/p\pmod p\quad(0\le j<p),
\quad U=u_0,\quad V=P_{\ell-1}(\alpha)^p.
\]

By \eqref{eq:old5.2}, all these quotients are integral. The recurrence identities in this step are understood in the residue field; here \(\alpha\) and \(P_j(\alpha)\) also denote their reductions. The adjacent polynomials \(\overline P_{\ell},\overline P_{\ell-1}\) have no common root. Otherwise, repeatedly lowering the index by the three-term recurrence would give \(P_0=0\). All integers by which we must divide in this recurrence are less than \(p\), so division remains valid after reduction. This proves that \(V\ne0\). From \eqref{eq:old1.4}, \eqref{eq:old5.2}, and \(P_{p-1}\equiv1\), we obtain

\[
u_1=\alpha U-\ell V,
\]

\[
(j+1)u_{j+1}=(2j+1)\alpha u_j-ju_{j-1}
\quad(1\le j\le p-2).
\]

The purpose of comparing these two sequences is to use their common recurrence to obtain a nonzero quantity independent of the index. The sequences \(P_j(\alpha)\) and \(u_j\) satisfy the same second-order recurrence and differ only in their initial values. Define
\[
\Delta_j=j\bigl(P_j(\alpha)u_{j-1}-P_{j-1}(\alpha)u_j\bigr).
\]
Since \(P_0=1\), \(P_1(\alpha)=\alpha\), and \(u_1=\alpha U-\ell V\), we have \(\Delta_1=\ell V\). Substituting the two recurrences into \(\Delta_{j+1}\) gives
\begin{align*}
\Delta_{j+1}
&=((2j+1)\alpha P_j-jP_{j-1})u_j
 -P_j((2j+1)\alpha u_j-ju_{j-1})\\
&=j(P_ju_{j-1}-P_{j-1}u_j)=\Delta_j,
\end{align*}
where \(P_j\) in this calculation means \(P_j(\alpha)\). Thus this difference of products equals the same nonzero constant for every \(j\):

\begin{equation}
j\bigl(P_j(\alpha)u_{j-1}-P_{j-1}(\alpha)u_j\bigr)=\ell V\ne0
\quad(1\le j<p).
\label{eq:old5.4}
\end{equation}

This nonzero constant will show that \(P_n(\alpha)\) and \(P_n'(\alpha)\) cannot both be divisible by excessively high powers of \(p\). The root of the reduction under consideration has multiplicity at least \(p\), so \(P_n'(\alpha)/p\) belongs to the valuation ring of \(K\). Denote its reduction by \(D\). When \(r\ge1\), dividing the derivative formula for \(P_n\) by \(p\) and reducing gives
\[
(1-\alpha^2)D=r(u_{r-1}-\alpha u_r).
\]
The other derivative formula gives
\[
(1-\alpha^2)P_r'(\alpha)
=r(P_{r-1}(\alpha)-\alpha P_r(\alpha)).
\]
Multiply these identities by \(P_r(\alpha)\) and \(u_r\), respectively, and subtract. The two terms containing \(\alpha P_r(\alpha)u_r\) cancel, and \eqref{eq:old5.4} yields
\begin{equation}
(1-\alpha^2)\bigl(P_r(\alpha)D-P_r'(\alpha)u_r\bigr)=\ell V\ne0.
\label{eq:old5.5}
\end{equation}
When \(r=0\), the relations \(n=\ell p\) and \(P_{\ell p-1}(\alpha)\equiv V\) directly give \((1-\alpha^2)D=\ell V\), so the same identity remains valid.

It follows that if the reduction of \(P_r(\alpha)\) is zero, then \(u_r\ne0\); that is, \(P_n(\alpha)\) is divisible by \(p\) exactly once. If \(u_r=0\), then the reduction of \(P_r(\alpha)\) is nonzero and \(D\ne0\); that is, \(P_n(\alpha)\) is divisible by \(p^2\) or is zero, whereas \(P_n'(\alpha)\) is divisible by \(p\) exactly once. We shall next use this information about the constant and linear coefficients to determine which factorizations are possible.

\textbf{(iv) How roots with the same reduction are distributed between the two factors.}

Put \(x=\alpha+Y\). The roots reducing to \(\overline\alpha\) are exactly those having positive valuation in the variable \(Y\). Let \(E(Y)\) be the monic polynomial formed by these roots. Hensel factorization, or the Weierstrass preparation theorem, gives
\[
P_n(\alpha+Y)=E(Y)J(Y),\qquad J(Y)\in\mathcal O_K[[Y]]^\times.
\]
Here \(\mathcal O_K\) is the valuation ring of \(K\), and \(J(Y)\) is a power series with unit constant term that contains the factors associated with the other roots of the reduction. The degree of \(E\) is the multiplicity of the root of the reduction under consideration, namely \(p\) or \(p+1\). Moreover, \(E(Y)\equiv Y^{\deg E}\pmod p\), so every coefficient other than the leading coefficient is divisible by \(p\).

The preceding step determined the constant and linear coefficients of the entire polynomial \(P_n(\alpha+Y)\). Since \(J(0)\) is a unit,
\[
v_p(E(0))=v_p(P_n(\alpha)).
\]
If this valuation is at least \(2\), or if both constant terms are zero, then the second term in
\[
P_n'(\alpha)=E'(0)J(0)+E(0)J'(0)
\]
is divisible by \(p^2\). Consequently, \(v_p(P_n'(\alpha))=1\) also implies \(v_p(E'(0))=1\). The required information about the constant and linear coefficients has therefore been transferred to \(E\). We now consider two cases.

If \(u_r\ne0\), the factor has degree \(p\) or \(p+1\), its constant term has valuation \(1\), and every other coefficient except the leading coefficient is divisible by \(p\). These are precisely the Eisenstein conditions: the leading coefficient is a unit, all the other coefficients are divisible by \(p\), and the constant term is not divisible by \(p^2\). Hence \(E\) is irreducible over \(K\). Since the coefficients of \(F,G\) also lie in \(K\), they cannot each contain only some of the roots of \(E\). The entire group must belong to the same factor.

If \(u_r=0\), the degree is \(p\), and the linear coefficient has valuation exactly \(1\). If \(h=v_p(P_n(\alpha))\ge2\) is finite, the vertices of the Newton polygon are

\[
(0,h),\quad(1,1),\quad(p,0).
\]

Thus this group of roots consists of a linear factor and an Eisenstein factor of degree \(p-1\). In the \(Y\)-coordinate, the root of the linear factor has valuation \(h-1\ge1\), whereas the remaining roots have valuation \(1/(p-1)\). If \(P_n(\alpha)=0\), we remove the factor \(Y\) directly; the remaining factor is again Eisenstein of degree \(p-1\). The same local factorization therefore holds in this case as well.

In the second case, denote the root of the linear factor by \(y_0\) and the remaining \(p-1\) roots by \(y_1,\ldots,y_{p-1}\). The Newton polygon gives
\[
v_p(y_0)\ge1,\qquad v_p(y_j)=\frac1{p-1}\quad(1\le j\le p-1),
\]
including the case \(y_0=0\). Since \(p\) is odd, \(1/(p-1)<1\), so the two valuations are different and
\[
v_p(y_0-y_j)=\frac1{p-1}.
\]
Returning from \(Y\) to the original variable \(x=\alpha+Y\) does not change root differences. If both irreducible factors belong to the same rational factor, there are no root differences between \(F\) and \(G\) to count. If they belong to \(F\) and \(G\), respectively, there are exactly \(p-1\) such root pairs. The total contribution of this group to the valuation of the resultant is then exactly
\[
(p-1)\cdot\frac1{p-1}=1.
\]
We have thus proved that the group associated with each nonzero multiple root of the reduction contributes at most \(1\) to \(v_p\operatorname{Res}_x(F,G)\).

Such roots of the reduction can only come from \(\overline P_{\ell}\). Its roots are all simple, and parity shows that it has a zero root exactly when \(\ell\) is odd. Hence the number of nonzero roots is \(\ell-(\ell\bmod2)=2\lfloor \ell/2\rfloor\). Differences between roots with distinct reductions have valuation zero; a simple root of the reduction cannot belong to both factors; and zero and infinity were excluded in step (ii). Summing the contributions therefore gives
\[
v_p\operatorname{Res}_x(F,G)\le2\lfloor \ell/2\rfloor.
\]
Finally, use \(\operatorname{Res}_x(F,G)=R^2\) to divide the left-hand side by \(2\). Since \(n=\ell p+r\), with \(0\le r<p\), we obtain
\[
v_p(R)\le\lfloor \ell/2\rfloor
=\lfloor n/(2p)\rfloor=\lfloor m/p\rfloor.
\]
This proves \eqref{eq:old5.1}.
\end{proof}

\subsection{From local estimates to a global upper bound}

We now sum the exponents over all odd primes. For a large prime, we know that \(v_p(R)\) is bounded by both \(2k-1\) and \(\lfloor m/p\rfloor\). If a prime occurs three times, count its \(\log p\) once at each of levels one, two, and three, instead of counting it as \(3\log p\) all at once. In general, a prime counted at level \(j\) must satisfy \(\lfloor m/p\rfloor\ge j\), or equivalently \(p\le m/j\). Thus the total at level \(j\) can be bounded using all primes not exceeding \(m/j\). The harmonic sum below arises when these levels are added together.

\begin{proof}[Proof of Proposition~\ref{prop:upper}]
Put \(T=\sqrt{2n}\). For odd primes \(p\le T\), Lemma~\ref{prop:uniform} shows that each prime contributes strictly less than \(k\log(2n)\). There are at most \(T\) such primes, so
\begin{equation}
\sum_{p\le T,\ p\text{ odd}}v_p(R)\log p
<kT\log(2n).
\label{eq:old7.1}
\end{equation}
For \(p>T\), we have \(\log_p(2n)<2\). Lemma~\ref{prop:uniform}, the integrality of the valuation, and Lemma~\ref{prop:large} give
\[
v_p(R)\le\min\left(2k-1,\left\lfloor\frac mp\right\rfloor\right).
\]
Define \(\vartheta(y)=\sum_{p\le y}\log p\), where the sum runs over all primes. Counting by levels yields
\begin{align*}
\sum_{p>T}v_p(R)\log p
&\le\sum_{p>T}\sum_{j=1}^{2k-1}\mathbf1_{\{p\le m/j\}}\log p\\
&=\sum_{j=1}^{2k-1}\sum_{T<p\le m/j}\log p
\le\sum_{j=1}^{2k-1}\vartheta(m/j).
\end{align*}
Here \(\mathbf1_{\{p\le m/j\}}\) is \(1\) when the condition holds and \(0\) otherwise. The elementary bound for the product of primes is
\begin{equation}
\vartheta(x)<2x\log2\quad(x\ge1)
\label{eq:old7.2}
\end{equation}
To prove it, proceed by induction on integers. The product of the primes in \((u,2u]\) divides \(\binom{2u}{u}\le2^{2u-1}\); the product of the primes in \((u+1,2u+1]\) divides \(\binom{2u+1}{u}\le2^{2u}\). Combining these inequalities with the induction hypotheses for \(u\) and \(u+1\), respectively, gives the even and odd cases. The cases \(1,2\) are checked directly. For real arguments, take the integer part. Consequently,
\[
\sum_{p>T}v_p(R)\log p
<2m\log2\sum_{j=1}^{2k-1}\frac1j
=2m\log2\,h_{2k-1}.
\]
Adding this to the contribution from small primes gives the upper bound in the proposition.
\end{proof}

\section{A lower bound for the odd part of the resultant}\label{sec:lower}

In this section we show that the same integer \(R_{\mathrm{odd}}\) cannot be too small. We now use the real absolute value: all roots of \(P_n\) lie in \((-1,1)\), and all roots of \(H\), expressed in the squared variable, lie in \((0,1)\).

The argument has three steps. First, in the identity \eqref{eq:root_identity} at a root, we already have an upper bound for the absolute value of \(W_{n-1}\), and hence obtain a lower bound for the absolute value of the derivative. Second, we multiply the derivatives at all roots belonging to the factor \(A\). This product contains the differences between roots of the two distinct factors \(A\) and \(B\), and therefore contains the resultant we wish to estimate. It also contains differences between roots of \(A\) itself. Third, we divide out these additional factors and give an upper bound for this divisor. We first state the resulting estimate and then carry out these three steps.

\begin{proposition}\label{prop:lower}
With the notation of Section~\ref{sec:construction}, we have

\begin{equation}
\log R_{\mathrm{odd}}
>2km\log2-k\big[(s(e)+5)\log2+\log h_n+\log k\big]
+2\log2.
\label{eq:old6.5}
\end{equation}

When \(k=1\), we also have

\begin{equation}
\log R_{\mathrm{odd}}
>m\log6-(s(m)+1)\log2-\log3-\log h_n.
\label{eq:old6.6}
\end{equation}
\end{proposition}

\subsection{From derivatives at roots to the resultant}

As before, let \(t_1,\ldots,t_k\) be the roots of \(A\), and define
\[
\mathcal V_A=\prod_{1\le i<j\le k}(t_i-t_j)^2,
\]
where the empty product is taken to be \(1\) when \(k=1\). This quantity contains only differences between roots of \(A\), and involves no roots of \(B\).

Differentiating \(H=AB\) gives \(H'=A'B+AB'\). Substituting \(A(t_i)=0\) makes the second term vanish, so that
\[
H'(t_i)=A'(t_i)B(t_i).
\]
Here \(B(t_i)\) contains the differences between \(t_i\) and the roots of \(B\), which are precisely the factors needed for the resultant. The factor \(A'(t_i)\) contains the differences between \(t_i\) and the other roots of \(A\), which are the additional factors introduced by this procedure. On multiplying over all \(i\), each pair of distinct roots of \(A\) occurs twice. After taking absolute values, these factors give \(\mathcal V_A\) as defined above, together with a power of the leading coefficient. We now write this relation precisely.

\begin{lemma}\label{lem:real_product}
\begin{equation}
|H'(t)|\ge
\frac{\kappa_n}{2h_n t^{(\varepsilon+1)/2}(1-t)}
>\frac{\kappa_n}{2h_n(1-t)}.
\label{eq:old6.1}
\end{equation}
Here \(t\) is any root of \(H\). Moreover, the resultant satisfies
\begin{equation}
|R|=a^{m-2k}\frac{\prod_i|H'(t_i)|}{\mathcal V_A}
>\frac{a^{m-2k+1}}{A(1)\mathcal V_A}
\left(\frac{\kappa_n}{2h_n}\right)^k.
\label{eq:old6.2}
\end{equation}
\end{lemma}

\begin{proof}
If \(P_n(x)=0\) and \(-1<x<1\), then \eqref{eq:root_identity} and \(|W_{n-1}(x)|\le h_n\) give
\[
|P_n'(x)|\ge\frac1{h_n(1-x^2)}.
\]
Differentiating \(\kappa_nP_n(x)=x^\varepsilon H(x^2)\) at a nonzero root gives
\[
\kappa_n P_n'(x)=2x^{\varepsilon+1}H'(x^2).
\]
Setting \(t=x^2\in(0,1)\) proves \eqref{eq:old6.1}. On the other hand,
\[
\prod_{i=1}^k|A'(t_i)|=a^k\mathcal V_A,\qquad
|R|=a^e\prod_{i=1}^k|B(t_i)|.
\]
Multiplying \(H'(t_i)=A'(t_i)B(t_i)\) over all \(i\), we obtain
\[
|R|=a^{e-k}\frac{\prod_i|H'(t_i)|}{\mathcal V_A}.
\]
Using \(e-k=m-2k\), \eqref{eq:old6.1}, and
\(\prod_i(1-t_i)=A(1)/a\), we obtain \eqref{eq:old6.2}.
\end{proof}

\subsection{Controlling differences between roots of the same factor}

In the preceding lemma, \(\mathcal V_A\) occurs in the denominator. To show that the whole quotient is sufficiently large, we need to show that this denominator is sufficiently small. Although each root difference has absolute value less than \(1\), bounding them individually by \(1\) would not give a lower bound strong enough for the final comparison.

The Vandermonde determinant is exactly the product of the differences over all pairs of distinct roots, so it allows us to estimate these differences together. We use Chebyshev polynomials because their absolute values on the interval are at most \(1\), while their leading coefficients can be calculated explicitly. Replacing the monomials in the determinant by these polynomials makes every matrix entry bounded in absolute value and introduces an explicit power of two through the change of basis. This power of two supplies the improvement we need.

\begin{lemma}\label{lem:vandermonde}
For any \(k\) points in \([0,1]\), we have
\begin{equation}
\mathcal V_A\le\frac{k^k}{2^{2(k-1)^2}}.
\label{eq:old6.4}
\end{equation}
\end{lemma}

\begin{proof}
Let \(T_j\) denote the Chebyshev polynomial of the first kind, defined by
\[
T_j(\cos\theta)=\cos(j\theta).
\]
Thus \(|T_j(2t-1)|\le1\) for \(t\in[0,1]\). Consider the evaluation matrix of
\(1,T_1(2t-1),\ldots,T_{k-1}(2t-1)\) at the given points. The Euclidean norm of every column is at most \(\sqrt k\), so Hadamard's inequality bounds the square of its determinant by \(k^k\).

On the other hand, the leading coefficient of \(T_j(2t-1)\) is \(2^{2j-1}\) for \(j\ge1\). The change-of-basis matrix from the monomial basis to this basis is triangular, with determinant
\[
\prod_{j=1}^{k-1}2^{2j-1}=2^{(k-1)^2}.
\]
The square of the original Vandermonde determinant is \(\mathcal V_A\). Therefore
\[
2^{2(k-1)^2}\mathcal V_A\le k^k,
\]
which proves the assertion. When \(k=1\), both sides of the claimed inequality equal \(1\).
\end{proof}

After taking logarithms, this estimate contributes \(2(k-1)^2\log2\). It improves the lower bound, whose previous leading term depended on the product \(ke\) of the two factor degrees, to one with leading term \(2km\log2\). This allows a uniform comparison with the upper bound in Section~\ref{sec:upper}.

\subsection{Combining the lower bounds and improving the case of a linear factor}

\begin{proof}[Proof of Proposition~\ref{prop:lower}]
By Lemma~\ref{lem:endpoint}, \(a>A(1)=2^{2k-s(k)}\), and \(m-2k\ge0\). Substituting this into \eqref{eq:old6.2}, dividing out \(2^{v_2(R)}\), and using \eqref{eq:old2.8}, we obtain
\[
R_{\mathrm{odd}}>
\frac{2^{(2k-s(k))(m-2k)+k(n-s(n)-1)-(2k-s(k))e}}
     {h_n^k\mathcal V_A}.
\]
Since \(n-s(n)=2m-s(m)\) and \(s(m)=s(k)+s(e)\), the exponent in the numerator simplifies to \(k(2e-s(e)-1)\). Thus
\begin{equation}
\log R_{\mathrm{odd}}
>k[2e-s(e)-1]\log2-k\log h_n-\log\mathcal V_A.
\label{eq:old6.3}
\end{equation}
By Lemma~\ref{lem:vandermonde},
\[
-\log\mathcal V_A\ge2(k-1)^2\log2-k\log k.
\]
Expanding with \(e=m-k\) proves \eqref{eq:old6.5}.

When \(k=1\), the general lower bound loses some additional information about the leading coefficient. In this case \(A(1)=2\), and \(a\) is a positive odd integer greater than \(2\), so \(a\ge3\). The restriction on degrees from their binary expansions also implies that \(m\) is odd, and \eqref{eq:old2.7} gives \(v_2(R)=m-1\). Retaining \(a\ge3\) in \eqref{eq:old6.2} and using \(\mathcal V_A=1\), we obtain
\[
R_{\mathrm{odd}}>
\frac{3^{m-1}2^{m-s(m)-1}}{h_n}.
\]
Taking logarithms gives \eqref{eq:old6.6}.
\end{proof}

\section{Comparison of the bounds and proof of the main theorem}\label{sec:conclusion}

We have now established both estimates for the same positive integer \(R_{\mathrm{odd}}\). It remains to compare them. For a factor of degree \(k\ge2\), after dividing both bounds by \(k\), the coefficient of the leading term is \(2\) in the lower bound and at most \(11/6\) in the upper bound. The difference between these terms grows linearly and eventually exceeds the error terms in the estimates. For a linear factor, we use the improvement supplied by the inequality \(a\ge3\) retained in the preceding section.

\begin{proposition}\label{prop:asymptotic}
For every integer \(n\ge100000\), the polynomial \(M_n\) is irreducible in \(\mathbb Q[x]\).
\end{proposition}

\begin{proof}
Suppose that a factorization as in Section~\ref{sec:construction} exists. First, we express the error terms in terms of \(n\) alone. The bound on the number of binary digits, together with
\[
h_n\le1+\int_1^n\frac{dt}{t}=1+\log n,
\]
gives, for \(n\ge64\),
\begin{equation}
(s(e)+5)\log2+\log h_n+\log k
\le2\log n+6\log2+\log(1+\log n)
\le3\log(2n).
\label{eq:old7.4}
\end{equation}
The last inequality is equivalent to \(8(1+\log n)\le n\). It holds at \(n=64\), and the difference between its right- and left-hand sides is strictly increasing thereafter.

\textbf{(i) \(k\ge2\).} The harmonic sum satisfies
\begin{equation}
h_{2k-1}\le\frac{11}{12}k.
\label{eq:old7.5}
\end{equation}
Equality holds when \(k=2\). Increasing \(k\) to \(k+1\) adds only \(1/(2k)+1/(2k+1)<11/12\) to the left-hand side, so the assertion follows by induction.

Set \(X=\log R_{\mathrm{odd}}/k\). Proposition~\ref{prop:lower} and \eqref{eq:old7.4}, after discarding a positive term, give
\begin{equation}\label{eq:comparison_lower}
X>2m\log2-3\log(2n).
\end{equation}
Proposition~\ref{prop:upper} and \eqref{eq:old7.5} give
\begin{equation}\label{eq:comparison_upper}
X<\frac{11}{6}m\log2+\sqrt{2n}\log(2n).
\end{equation}
The lower bound minus the upper bound is exactly
\begin{equation}\label{eq:comparison_gap}
\frac{m\log2}{6}-(\sqrt{2n}+3)\log(2n).
\end{equation}
If this difference is positive, the same \(X\) must exceed the larger number and be less than the smaller number, which is a contradiction. Since \(m\ge(n-1)/2\), a sufficient condition is
\begin{equation}
\frac{(n-1)\log2}{12}>
(\sqrt{2n}+3)\log(2n),
\label{eq:old7.6}
\end{equation}
\textbf{(ii) \(k=1\).} In this case \(m\) is odd and \(s(m)=s(e)+1\). Write
\[
E_n=(s(m)+1)\log2+\log3+\log h_n.
\]
Since \(\log3<3\log2\), inequality \eqref{eq:old7.4} gives \(E_n<3\log(2n)\). The lower and upper bounds are, respectively,
\[
\log R_{\mathrm{odd}}>m\log6-E_n,
\qquad
\log R_{\mathrm{odd}}<2m\log2+\sqrt{2n}\log(2n).
\]
The difference between the leading terms is \(m\log(3/2)>m\log2/6\). Thus the same sufficient condition \eqref{eq:old7.6} also excludes a linear factor.

\textbf{(iii) Verifying the degree threshold.}
At \(n_0=100000\), the left-hand side is greater than

\[
99999\cdot0.69/12=5749.9425,
\]

and the right-hand side is less than

\[
(448+3)\cdot12.3=5547.3.
\]

These estimates follow directly from the following rational inequalities:

\[
\log2>2\left(\frac13+\frac1{81}+\frac1{1215}\right)
=\frac{842}{1215}>0.69;
\]

\[
\exp(3)>\sum_{j=0}^8\frac{3^j}{j!}
=\frac{89641}{4480}>20,\qquad \exp(0.3)>1.3.
\]

Hence \(\exp(12.3)>20^4\cdot1.3=208000>200000\), so \(\log200000<12.3\); also, \(\sqrt{200000}<448\).

After dividing \eqref{eq:old7.6} by \(n\), the left-hand side \((1-1/n)\log2/12\) is strictly increasing, whereas the right-hand side

\[
\frac{\sqrt2\log(2n)}{\sqrt n}+\frac{3\log(2n)}n
\]

is strictly decreasing for \(n\ge100000\). Thus \eqref{eq:old7.6} holds for all such \(n\).
Consequently, for every \(n\ge100000\), any nontrivial factorization would lead to the contradiction above. Therefore \(M_n\) is irreducible.
\end{proof}

\begin{proof}[Proof of Theorem~\ref{thm:main}]
Proposition~\ref{prop:asymptotic} treats all \(n\ge100000\). For \(2\le n<100000\), we apply the computational result of Groth~\cite{groth}: \(P_n(x)\) for even \(n\), and \(P_n(x)/x\) for odd \(n\), are irreducible in \(\mathbb Q[x]\) for every original degree \(2\le n\le10^7\). The cited range contains the required finite interval by a wide margin.

Finally, \(M_n=\kappa_nP_n/x^\varepsilon\) differs from \(P_n/x^\varepsilon\) only by the nonzero rational factor \(\kappa_n\), so the two polynomials have the same irreducibility property. The theorem therefore holds for every \(n\ge2\).
\end{proof}

\begin{thebibliography}{9}
\bibitem{stieltjes} T. J. Stieltjes, Letter no. 275 to C. Hermite, 2 October 1890, in B. Baillaud and H. Bourget (eds.), \emph{Correspondance d'Hermite et de Stieltjes}, tome II, Gauthier-Villars, Paris, 1905, pp. 104--106.
\bibitem{wahab} J. H. Wahab, New cases of irreducibility for Legendre polynomials, \emph{Duke Math. J.} \textbf{19} (1952), no. 1, 165--176. \url{https://doi.org/10.1215/S0012-7094-52-01917-0}.
\bibitem{cullinan} J. Cullinan and F. Hajir, On the Galois groups of Legendre polynomials, \emph{Indag. Math. (N.S.)} \textbf{25} (2014), no. 3, 534--552. \url{https://doi.org/10.1016/j.indag.2014.01.004}.
\bibitem{groth} R. S. Groth, \emph{Generalized Sierpi\'nski Numbers and New Irreducibility Criteria on Legendre Polynomials}, Ph.D. dissertation, University of South Carolina, 2025. \url{https://scholarcommons.sc.edu/etd/8101/}.
\bibitem{noe} T. D. Noe, On the divisibility of generalized central trinomial coefficients, \emph{J. Integer Seq.} \textbf{9} (2006), Article 06.2.7. \url{https://cs.uwaterloo.ca/journals/JIS/VOL9/Noe/noe35.html}.
\bibitem{dlmf} NIST Digital Library of Mathematical Functions, Chapters 14 and 18, Equations (14.7.4), (14.7.6), (14.7.7), (14.12.13), (18.5.1), (18.14.1). \url{https://dlmf.nist.gov/}.
\end{thebibliography}

\end{document}
